\documentclass[11pt]{article}
\usepackage[utf8]{inputenc}
\usepackage{amsmath,amssymb,amsthm}
\usepackage[margin=1in]{geometry}
\usepackage{hyperref}
\usepackage{xcolor}
\usepackage{enumitem}
\usepackage{newunicodechar}
\newunicodechar{ℝ}{\ensuremath{\mathbb{R}}}
\newunicodechar{ℕ}{\ensuremath{\mathbb{N}}}
\newunicodechar{σ}{\ensuremath{\sigma}}
\newunicodechar{β}{\ensuremath{\beta}}
\newunicodechar{φ}{\ensuremath{\varphi}}
\newunicodechar{ψ}{\ensuremath{\psi}}
\newunicodechar{λ}{\ensuremath{\lambda}}
\newunicodechar{α}{\ensuremath{\alpha}}
\newunicodechar{ℓ}{\ensuremath{\ell}}
\newunicodechar{∑}{\ensuremath{\sum}}
\newunicodechar{∫}{\ensuremath{\int}}
\newunicodechar{∂}{\ensuremath{\partial}}
\newunicodechar{∈}{\ensuremath{\in}}
\newunicodechar{≠}{\ensuremath{\neq}}
\newunicodechar{≤}{\ensuremath{\leq}}
\newunicodechar{·}{\ensuremath{\cdot}}
\newunicodechar{²}{\ensuremath{^2}}
\newunicodechar{ᵐ}{\ensuremath{^{\mathrm{m}}}}
\newunicodechar{↦}{\ensuremath{\mapsto}}

\newcommand{\GR}{\mathrm{GibbsRegularity}}
\newcommand{\GO}{\mathrm{GibbsObservable}}

\newcommand{\userbox}[1]{\par\medskip\begin{quote}\small\textbf{\textcolor{blue!60}{DM:}} #1\end{quote}\medskip}
\newcommand{\claudebox}[1]{\par\medskip\begin{quote}\small\textbf{\textcolor{orange!60!black}{Claude:}} #1\end{quote}\medskip}
\newcommand{\gptbox}[1]{\par\medskip\begin{quote}\small\textbf{\textcolor{green!50!black}{GPT-5.5 Pro:}} #1\end{quote}\medskip}

\emergencystretch=5em

\title{Tide retrospective:\\\emph{Harmonic \texorpdfstring{$\GO$}{GibbsObservable}
  instances for monomials}}
\author{Daniel Murfet (with Claude Opus 4.7)}
\date{7 May 2026}

\begin{document}
\maketitle

\begin{abstract}
A single Tide-loop excursion (\textbf{G4} from the 7 May
candidates survey) closing the three carried
\texttt{Threepoint.GibbsObservable} hypotheses on Tide 13's
harmonic-cross-susceptibility-derivative theorem. Headline:
\texttt{Threepoint.harmonic\_id\_gibbsObservable\_pow} for any
\(k : \mathbb{N}\), proved by dominated differentiation under the
integral with a Gaussian-shaped dominator from Young's inequality.
397 lines of Lean; zero \texttt{sorry}, \texttt{axiom}, \texttt{native\_decide};
clean diagnostics. The session's load-bearing event was a GPT-5.5~Pro
consult that pivoted the dominator-integrability argument from a
piecewise \(|x| \le 1\) / \(|x| > 1\) split to a single global rpow-cast
route, and supplied the \texttt{nlinarith} hints (\((\lambda x \pm 2h)^2 \ge 0\))
for Young after a first attempt accumulated seven distinct frictions
across a 260-line draft.
\end{abstract}

\section{Setting}

This Tide is the natural close-out of Tide~13, which proved
\(\partial_h \mathrm{Cov}_h(x, x)\big|_{h=0} = 0\) for the harmonic Gibbs
measure but had to carry three \(\GO\) hypotheses (for \(x\), \(x^2\),
\(x^3\)) as parameters because the harmonic monomial \(\GO\) instances
weren't yet in the seabed. G4 supplies them, generic in
\(k : \mathbb{N}\), making the Tide~13 theorem unconditional for the
canonical monomial cases.

The mathematical content is mild: for the harmonic potential
\(L(x) = \lambda x^2/2\) with linear perturbation \(A(x) = x\), the
perturbed numerator
\(N_k(h) := \int x^k \exp(-t(L(x) + h A(x)))\,dx\)
is differentiable in \(h\) at \(h = 0\) with derivative
\(\int x^k \cdot (-t \cdot x) \cdot \exp(-t L(x))\,dx\). This is dominated
differentiation under the integral, valid because the perturbed
integrand has a uniform Gaussian-shaped dominator on
\(|h| \le 1\)\footnote{Young's inequality: \(|h \cdot x| \le (\lambda/4) x^2 + h^2/\lambda\)
absorbs the linear perturbation into the quadratic decay.}.

The Step~2 deliberation (Claude $\leftrightarrow$ GPT-5.5~Pro) drafted
five candidates (a finite-\(k\) closed-form route, a generic-\(k\)
closed-form route, the dominated-differentiation route, a global-Cov
constancy result, and a `\textit{just close G4}'-shaped variant). Both
parties backed the dominated-differentiation route (Candidate D) — generic
in \(k\) at no extra cost, no Gaussian-moment library required.

\section{The thing formalised}

\begin{quote}\itshape
For \(\lambda > 0\), \(t > 0\), and any \(k \in \mathbb{N}\), the
canonical monomial \((x \mapsto x^k)\) is a \(\GO\) for the harmonic
potential \((\lambda/2) x^2\) with linear perturbation \(A(x) = x\)
under Lebesgue measure on \(\mathbb{R}\): the perturbed numerator
\(\int x^k \cdot \exp(-t(L(x) + h A(x)))\,dx\) reduces to the
unperturbed integral at \(h = 0\), and is differentiable at \(h = 0\)
with the expected pointwise-derivative integral.
\end{quote}

\noindent The Lean signature:
\begin{verbatim}
theorem Threepoint.harmonic_id_gibbsObservable_pow
    {lam t : ℝ} (hlam : 0 < lam) (ht : 0 < t) (k : ℕ) :
    Threepoint.GibbsObservable (volume : Measure ℝ)
      (fun x : ℝ => lam / 2 * x ^ 2)
      (fun x : ℝ => x) t (fun x : ℝ => x ^ k)
\end{verbatim}

\noindent located in
\texttt{Laplace/OneD/HarmonicGibbsObservableMonomials.lean}, namespace
\texttt{Laplace.OneD}.

The result is the natural pair to Tide~10's harmonic
\(\GR\) instance: where the \(\GR\) instance handles the
\emph{partition-function} differentiability uniformly in \(\varphi\)
(via the closed-form route), this \(\GO\) instance handles the
\emph{numerator} differentiability per-observable. Together they
discharge the analytic content of the harmonic FDT and
cross-susceptibility theorems, instantiated at any monomial.

\section{Proof strategy}

The Lean proof is built from five small lemmas, each with a clean
mathematical narrative:

\paragraph{Young's inequality with weight \(\lambda/2\).} For
\(\lambda > 0\), \(|h \cdot x| \le (\lambda/4) x^2 + h^2/\lambda\). The
two one-sided bounds come from \((\lambda x - 2h)^2 \ge 0\) and
\((\lambda x + 2h)^2 \ge 0\) after dividing through by \(4 \lambda\).
The Lean idiom that landed cleanly: clear the \(/\lambda\) by
\texttt{field\_simp} + \texttt{ring}, then \texttt{nlinarith} with the
square hint.

\paragraph{Dominator integrability.}
\(\int |x|^k \exp(-c x^2)\,dx < \infty\) for \(c > 0\), \(k \in \mathbb{N}\).
Lean route: get
\(\int x^{(k:\mathbb{R})} \exp(-c x^2)\,dx < \infty\)
from \texttt{integrable\_rpow\_mul\_exp\_neg\_mul\_sq} (Mathlib's
real-power Gaussian moment), cast to \texttt{Nat}-power via
\texttt{Real.rpow\_natCast}, then bridge \(x^k\) to \(|x|^k\) via
\texttt{Integrable.norm} (since \(\|x^k \cdot e^{-cx^2}\| = |x|^k \cdot e^{-cx^2}\)
because the exponential is positive).

\paragraph{Pointwise integrand bound.}
\(\|x^k \exp(-t(\lambda/2 \cdot x^2 + h x))\|
\le e^{t/\lambda} \cdot |x|^k \cdot e^{-(t\lambda/4) x^2}\)
for \(|h| \le 1\). Direct application of Young followed by
\(\exp\)-monotonicity.

\paragraph{Pointwise differentiability.} For each fixed \(x\), the
function
\(h \mapsto x^k \cdot \exp(-t(\lambda/2 \cdot x^2 + h x))\)
is differentiable in \(h\) by the standard
\texttt{HasDerivAt.exp.const\_mul} chain. The derivative is
\(x^k \cdot (-(t \cdot x)) \cdot \exp(-t(\lambda/2 \cdot x^2 + h x))\).

\paragraph{Headline assembly via dominated differentiation.} Apply the
parametric-integral derivative
lemma\footnote{\texttt{MeasureTheory.hasDerivAt\_integral\_of\_dominated\_loc\_of\_deriv\_le}.}
with the four hypothesis families:
(i) the local neighbourhood
\texttt{Metric.ball 0 1 ∈ nhds 0};
(ii) AE strong measurability via \texttt{fun\_prop};
(iii) integrability of the derivative integrand bounded by the
dominator at index \(k+1\) (the derivative integrand is
\(-t \cdot x^{k+1} \cdot \exp(-t(\ldots))\));
(iv) pointwise \texttt{HasDerivAt} from the chain rule above.
The result's second component is the desired \texttt{HasDerivAt}; a
small \texttt{integral\_congr\_ae} massages \(0 \cdot x = 0\) inside
the perturbed integrand to match the \(\GO\) shape.

\section{Roadblocks and resolutions}

\paragraph{The first attempt accumulated seven distinct errors.} I
drafted all the analytic-core lemmas in one pass (\(\sim\)260 lines)
and tried to compile. The build report showed three \texttt{nlinarith}
failures (Young's inequality, AM-GM bounds), two \texttt{Real.rpow}
type-mismatches (the dominator integrability), one \texttt{mul\_le\_mul}
argument-order issue, and one bound-shape mismatch. Each error was
individually fixable but the combined diagnosis cost would have been
slow. The lean-formalisation skill's \emph{<50\% probability of direct
proof \(\implies\) factor first} rule applied; I reverted and wrote
the handoff note while launching a focused GPT-5.5~Pro consult.

\paragraph{The GPT consult was decisive.} Three pieces of advice
landed:
\gptbox{For Young: don't ask \texttt{nlinarith} to reason directly with
\texttt{/ lam}. Multiply by \(4\lambda\) first, then \texttt{nlinarith}
with \((\lambda x \mp 2h)^2 \ge 0\).}
\gptbox{For dominator integrability: there's a \emph{global} dominator
\(M \cdot e^{-(c/2) x^2}\), \(M = e^{k^2/(2c)}\), via the same Young
applied to \((k, |x|)\). No piecewise \(|x| \le 1\) / \(|x| > 1\) split
needed.}
\gptbox{For \texttt{HasDerivAt}: the canonical chain is\\
\texttt{(((((const).add (id.mul\_const)).const\_mul).neg).exp).const\_mul}.}
With these in hand the second attempt landed clean.

In the end, the dominator integrability route I used differed from
GPT's advice: rather than the global Young-based AM-GM dominator,
I used \texttt{Integrable.norm} on
\texttt{integrable\_rpow\_mul\_exp\_neg\_mul\_sq} at \texttt{s = (k:ℝ)},
which is even simpler. (\(\sim\)10 lines vs \(\sim\)50 for the AM-GM
route.) Worth recording — both routes work; the Mathlib
\texttt{Integrable.norm} bridge between \(x^k\) and \(|x|^k\) is
ergonomically excellent for this case.

\paragraph{Pi.add unfolding edge-case.} The \texttt{HasDerivAt.add}
of \texttt{(fun \_ => c)} and \texttt{(fun y => y * x)} produces a
goal in \texttt{((fun \_ => c) + (fun y => y * x))} form (with
\texttt{Pi.add}), not the inline-lambda form
\texttt{(fun y => c + y * x)}. \texttt{simpa} normalised \(0 + x\)
inside the goal too aggressively. Resolved by introducing a manual
\texttt{hzero : (0 : ℝ) + x = x = zero\_add x} and \texttt{rwa}-ing
it into the \texttt{hsum} term \emph{before} using.

\paragraph{Continuity proofs via \texttt{fun\_prop}.} Manual continuity
chains\footnote{Of the form
\texttt{Continuous.const\_mul.add.const\_mul.neg.exp}.} hit the same
\texttt{Pi.add} boundary as the \texttt{HasDerivAt} issue. Switching
to \texttt{by fun\_prop} sidestepped this entirely. Worth a
\texttt{CLAUDE.md} note for the next dominated-differentiation Tide.

\section{What was Mathlib, what was new}

\paragraph{Mathlib provided.}
\texttt{MeasureTheory.hasDerivAt\_integral\_of\_dominated\_loc\_of\_deriv\_le}
(the main parametric-integral lemma);
\texttt{integrable\_rpow\_mul\_exp\_neg\_mul\_sq} (real-power Gaussian
moment integrability); the \texttt{Integrable.norm} /
\texttt{Real.rpow\_natCast} bridge; the \texttt{HasDerivAt.exp},
\texttt{.const\_mul}, \texttt{.add}, \texttt{.neg}, \texttt{.mul\_const}
combinators; \texttt{Real.exp\_le\_exp.mpr}, \texttt{Real.exp\_add},
\texttt{abs\_le}, \texttt{abs\_pow}, \texttt{Real.norm\_eq\_abs}.

\paragraph{The seabed provided.}
\texttt{Threepoint.GibbsObservable} (the target structure) and Tide
13's \texttt{cov\_h\_id\_id\_deriv\_harmonic\_eq\_zero} (the consumer
of this Tide's output). No prior 1D Gaussian moment infrastructure
was used directly; the proof routes around the seabed by going
straight through Mathlib's parametric-integral API.

\paragraph{What was new in this tide.}
The Young's-with-weight inequality (likely promotable to
\texttt{lean-common}); the dominator integrability lemma; the
pointwise integrand bound; the pointwise \texttt{HasDerivAt} for the
perturbed monomial integrand; the headline \(\GO\) instance.

\section{Lessons}

\begin{itemize}
\item \emph{Strategic GPT consults at the friction point pay off
asymmetrically.} The first attempt accumulated seven distinct errors
in a 260-line draft; the GPT consult turned them into a clean second
pass. The pattern matches what the lean-formalisation skill describes
as the \(\sim\)5-minute consult that saves a multi-hour detour.

\item \emph{\texttt{nlinarith} can't divide.} For inequalities with
\texttt{/ c}, multiply through first via
\texttt{field\_simp + ring}, then invoke \texttt{nlinarith} with the
right square hint. Generic enough to lift to \texttt{laplace/CLAUDE.md}.

\item \emph{\texttt{fun\_prop} for continuity beats manual chains.}
The continuity proofs were the second-most-fragile part of the first
attempt. \texttt{fun\_prop} dispatched all of them in one tactic call.

\item \emph{\texttt{Integrable.norm} is the right bridge between
\texttt{x\^{}k} and \texttt{|x|\^{}k}.} Cleaner than piecewise
\(|x| \le 1\) / \(|x| > 1\) splits when the inner integrability already
exists in Mathlib for \(x^k\). One more pattern for the API reference.
\end{itemize}

\section{Follow-ups}

\begin{itemize}
\item \emph{Promote \texttt{abs\_mul\_le\_quarter\_lambda\_sq\_add} to
\texttt{lean-common}.} Standard Young's-with-weight inequality; if a
second use materialises (e.g. in \texttt{singfluence-lean} or in a
future bounded-prior Tide), worth the promotion.

\item \emph{Package the dominated-differentiation infrastructure as a
reusable \(\GO\) factory.} The four hypothesis discharges
(integrability of \(F\) at \(h=0\), AE strong measurability,
integrand bound, pointwise differentiability) are generic for any
\(\varphi\) with a Gaussian-shaped dominator. A future Tide could
extract \texttt{Threepoint.gibbsObservable\_of\_dominated} so that
each new observable family doesn't re-run the whole hypothesis chain.

\item \emph{Tide~13 unconditionalisation.} With G4 closed, the next
natural Tide is to specialise Tide~13's
\texttt{cov\_h\_id\_id\_deriv\_harmonic\_eq\_zero} to the
unconditional case for \(\varphi(x) = x\) and \(B(x) = x\) by
substituting in the new instances. \(\sim\)20--30 lines on top of
this Tide.

\item \emph{Candidate E (deferred at deliberation).} Global
covariance constancy: for the harmonic Gibbs measure with linear
perturbation, \(\mathrm{Cov}_h(x, x)\) equals the unperturbed value
\(1/(\lambda t)\) for \emph{all} \(h\), not merely vanishing
derivative at \(0\). Different theorem statement; closed-form route
via Tide~10's \(Z(h)\); \(\sim\)100 lines.
\end{itemize}

\section*{Acknowledgements}

GPT-5.5~Pro contributed the deliberation-stage advisory at Step~2
(picking Candidate~D over the closed-form-binomial route), and a
mid-formalisation tactical consult that pivoted the analytic-core
proof from a friction-stuck draft to a clean
landing.\footnote{Tide log:
\texttt{sri/projects/patterning/tide-log/2026-05-07-tide-harmonic-gibbsobservable-monomials.md}.
Tactical consult on Lean frictions:
\texttt{sri/projects/patterning/tide-log/gpt55\_g4\_frictions\_v1.md}.}

\end{document}
